\chapter{Conclusion}
\label{chap:conclusion}

This thesis developed a reproducible framework for relating pointwise MAIA sensitivity to local OCT information in matched Usher syndrome visits. The framework links the functional and structural modalities through saved registration and quality-control decisions, three-nearest-B-scan mapping, approximately 100~$\mu$m footprint-aware sampling and hierarchical pooling to a point-level representation. It then evaluates prespecified models using participant-grouped nested validation rather than treating neighbouring retinal samples as independent cases.

The primary analysis preserves the all-eligible cohort by using a label-blind RETFound representation; the B0--B4 structural analysis is appropriately secondary and availability-limited. The location-plus-RETFound ridge was the best primary model (MAE 5.485~dB, 95\% CI 4.871--6.177; RMSE 6.986~dB; bias $-0.060$~dB; Bland--Altman limits $-13.754$ to 13.634~dB), improving on the location-only ridge (MAE 7.650~dB) and the training-mean null (9.346~dB). These values must be interpreted alongside the participant-level uncertainty, the device-floor error and the remaining limits of agreement, not in isolation.

The study shows that, within this internally evaluated cohort, spatially aligned local OCT information improves prediction of recorded pointwise sensitivity beyond simple non-imaging references. It does not demonstrate that OCT replaces microperimetry, that a named retinal layer drives function, or that the approach generalises beyond this cohort. Its contribution is an evidence-traceable spatial and statistical framework through which those stronger claims can be tested in future data.
